\documentclass[UTF8,a4paper,11pt]{ctexart}

\usepackage{amsmath,amssymb,mathtools}
\usepackage{booktabs}
\usepackage{geometry}
\usepackage{xcolor}
\usepackage{hyperref}
\usepackage{listings}
\usepackage{enumitem}
\usepackage{tabularx}

\geometry{margin=2.4cm}
\hypersetup{
  colorlinks=true,
  linkcolor=blue!60!black,
  urlcolor=blue!60!black,
  citecolor=blue!60!black
}

\lstdefinestyle{pycode}{
  language=Python,
  basicstyle=\ttfamily\small,
  keywordstyle=\color{blue!60!black},
  commentstyle=\color{green!40!black},
  stringstyle=\color{red!50!black},
  showstringspaces=false,
  columns=fullflexible,
  keepspaces=true,
  breaklines=true,
  frame=single,
  rulecolor=\color{black!20},
  backgroundcolor=\color{black!2},
  tabsize=4
}
\lstset{style=pycode}

\title{Assignment 5 Alignment 项目修复报告}
\author{Codex 修复记录}
\date{2026 年 6 月 5 日}

\begin{document}
\maketitle

\section{修复范围概览}

本次修复按“会阻断运行的问题优先、算法正确性其次、测试与工程清理收尾”的顺序执行。当前机器无 GPU，因此所有验证均限定为 CPU 可执行路径；依赖 vLLM 的 rollout/evaluation 路径只做入口约束与导入安全处理，不在本机强行运行。

\begin{table}[h]
\centering
\begin{tabularx}{\textwidth}{lX X}
\toprule
类别 & 主要文件 & 修复要点 \\
\midrule
SFT & \texttt{sft\_helper.py}, \texttt{sft\_experiment.py} & 稳定熵、训练循环、CPU 加载、vLLM 跳过 \\
GRPO / Dr.GRPO & \texttt{grpo\_helper.py}, \texttt{grpo\_experiment.py} & clip loss、mask 统计、constant normalization、optimizer 复用 \\
EI / GRPO 入口 & \texttt{expert\_iteration\_experiment.py}, \texttt{grpo\_experiment.py} & 无 GPU 时明确拒绝 vLLM rollout \\
数据与路径 & \texttt{data\_utils.py}, \texttt{paths.py} & 数据 schema 识别、路径校验延迟 \\
测试与 adapter & \texttt{tests/adapters.py}, \texttt{tests/test\_regressions.py} & adapter 透传、新增 CPU 回归测试 \\
\bottomrule
\end{tabularx}
\end{table}

\section{SFT：熵计算的数值稳定性}

\subsection{问题}

修复前的熵实现等价于先对 logits 做指数再归一化：

\begin{lstlisting}
# before
probs = torch.exp(logits) / torch.exp(logits).sum(dim=-1, keepdim=True)
entropy = -torch.sum(probs * torch.log(probs), dim=-1)
\end{lstlisting}

当 logits 很大时，例如 $[1000,999,998]$，$\exp(1000)$ 会 overflow，随后产生 \texttt{inf / inf} 或 \texttt{nan}。这会污染 SFT 日志，也会影响任何依赖 token entropy 的训练监控。

\subsection{数学推导}

令模型在位置 $t$ 的 logits 为 $\mathbf{z}_t \in \mathbb{R}^{V}$。真实概率为
\[
  p_i = \frac{\exp(z_i)}{\sum_{j=1}^{V}\exp(z_j)}.
\]
熵为
\[
  H(\mathbf{p}) = -\sum_{i=1}^{V} p_i \log p_i.
\]
直接计算 $\exp(z_i)$ 不稳定。使用 log-softmax：
\[
  \ell_i = \log p_i = z_i - \log\sum_{j=1}^{V}\exp(z_j).
\]
其中 $\log\sum_j \exp(z_j)$ 在 PyTorch 的 \texttt{log\_softmax} 内部使用 log-sum-exp 技巧：
\[
  \log\sum_j \exp(z_j)
  = m + \log\sum_j \exp(z_j-m),
  \quad m=\max_j z_j.
\]
此时 $z_j-m \le 0$，指数不会 overflow。最后
\[
  H = -\sum_i \exp(\ell_i)\ell_i.
\]

\subsection{修复后代码}

\begin{lstlisting}
# after: cs336_alignment/sft_helper.py
def compute_entropy(logits: torch.Tensor) -> torch.Tensor:
    log_probs = F.log_softmax(logits.float(), dim=-1)
    probs = log_probs.exp()
    return -torch.sum(probs * log_probs, dim=-1).to(logits.dtype)
\end{lstlisting}

这里先将 logits 转成 \texttt{float()}，避免半精度下 log-softmax 精度不足；返回时再转回原 dtype，保持调用方接口一致。

\section{SFT 损失、训练循环与尾部梯度累积}

\subsection{SFT 损失公式}

对于 prompt $x$ 和 response $y=(y_1,\dots,y_T)$，SFT 最大化 response token 的条件似然：
\[
  \max_\theta \sum_{t=1}^{T}\log \pi_\theta(y_t \mid x,y_{<t}).
\]
最小化形式为
\[
  \mathcal{L}_{\mathrm{SFT}}
  =
  -\frac{1}{B C}
  \sum_{i=1}^{B}\sum_{t=1}^{L}
  m_{i,t}\log\pi_\theta(y_{i,t}\mid x_i,y_{i,<t}),
\]
其中 $m_{i,t}\in\{0,1\}$ 是 response mask，$C$ 是 \texttt{normalize\_constant} 或 token 数相关归一化常数，$B$ 是 batch size。

\subsection{修复前问题}

训练循环曾将 \texttt{tqdm(dataloader)} 创建在 epoch 外层：

\begin{lstlisting}
# before
progress_bar = tqdm(dataloader, desc=f"Training ({len(train_dataset)} samples)")
for _ in range(train_config.epochs):
    for step, batch in enumerate(progress_bar):
        ...
\end{lstlisting}

这会导致 dataloader iterator 在第一轮 epoch 后耗尽。若 \texttt{epochs > 1}，后续 epoch 不再真正训练。

此外，尾部梯度累积窗口不足 \texttt{gradient\_accumulation} 时，日志仍除以固定累积步数，导致最后一个 optimizer step 的 loss/entropy/length 统计偏小。

\subsection{修复后代码}

\begin{lstlisting}
# after: sft_experiment.py / expert_iteration_experiment.py
for _ in range(train_config.epochs):
    progress_bar = tqdm(dataloader, desc=f"Training ({len(train_dataset)} samples)")
    for step, batch in enumerate(progress_bar):
        ...
        is_last_batch = step == len(dataloader) - 1
        if (step + 1) % train_config.gradient_accumulation == 0 or is_last_batch:
            ...
            accumulated_steps = (step % train_config.gradient_accumulation) + 1
            avg_loss = running_loss / accumulated_steps
            avg_entropy = running_entropy / accumulated_steps
            avg_len = running_len / accumulated_steps
\end{lstlisting}

修复后每个 epoch 都会重新创建 dataloader iterator；最后一个不足窗口也会执行 optimizer step，并按实际 microbatch 数统计日志。

\section{GRPO / Dr.GRPO：组内奖励归一化}

\subsection{公式}

GRPO 对同一个 prompt 采样 $G$ 个 response，得到 reward $r_{i,1},\dots,r_{i,G}$。组内 baseline 为
\[
  \mu_i = \frac{1}{G}\sum_{j=1}^{G}r_{i,j}.
\]
若启用标准差归一化，优势为
\[
  A_{i,j}
  =
  \frac{r_{i,j}-\mu_i}{\sigma_i+\epsilon}.
\]
本项目当前保留 PyTorch \texttt{std} 默认口径，即 sample standard deviation：
\[
  \sigma_i
  =
  \sqrt{\frac{1}{G-1}\sum_{j=1}^{G}(r_{i,j}-\mu_i)^2}.
\]
说明：GRPO 论文公式常见实现也可使用 population standard deviation，即分母为 $G$。本修复曾评估该口径，但原始 snapshot 测试锁定了 PyTorch 默认 sample std；为保持课程测试契约，本次保留 sample std，并在代码中新增 $G>1$ 防御。

\subsection{修复前问题}

当 $G=1$ 时，组内标准差没有统计意义；如果所有奖励相同，优势应为 0。缺少显式检查时，后续训练容易出现不清晰的 NaN 或无意义 advantage。

\begin{lstlisting}
# before
raw_rewards_per_group = raw_rewards.reshape((-1, group_size))
advantages = raw_rewards_per_group - raw_rewards_per_group.mean(-1, keepdim=True)
if normalize_by_std:
    advantages /= advantage_eps + raw_rewards_per_group.std(-1, keepdim=True)
\end{lstlisting}

\subsection{修复后代码}

\begin{lstlisting}
# after: cs336_alignment/grpo_helper.py
if group_size <= 1:
    raise ValueError("group_size must be greater than 1 for group-normalized rewards")

raw_rewards_per_group = raw_rewards.reshape((-1, group_size))
advantages = raw_rewards_per_group - raw_rewards_per_group.mean(-1, keepdim=True)
if normalize_by_std:
    advantages /= (advantage_eps + raw_rewards_per_group.std(-1, keepdim=True))
\end{lstlisting}

若组内 reward 全相同，则分子 $r_{i,j}-\mu_i=0$，因此 advantage 全为 0；新增回归测试覆盖了该性质。

\section{GRPO-Clip：裁剪目标与 clip fraction}

\subsection{公式}

令旧策略为 $\pi_{\theta_{\mathrm{old}}}$，当前策略为 $\pi_\theta$，单 token 概率比为
\[
  \rho_{i,t}(\theta)
  =
  \frac{\pi_\theta(y_{i,t}\mid x_i,y_{i,<t})}
       {\pi_{\theta_{\mathrm{old}}}(y_{i,t}\mid x_i,y_{i,<t})}
  =
  \exp\left(
    \log\pi_\theta(y_{i,t})-\log\pi_{\theta_{\mathrm{old}}}(y_{i,t})
  \right).
\]
GRPO/PPO 风格裁剪目标为
\[
  \mathcal{J}_{i,t}^{\mathrm{clip}}
  =
  \min\left(
    \rho_{i,t} A_i,
    \operatorname{clip}(\rho_{i,t},1-\epsilon,1+\epsilon)A_i
  \right).
\]
训练中最小化负目标：
\[
  \mathcal{L}_{i,t}^{\mathrm{clip}}
  =
  -\mathcal{J}_{i,t}^{\mathrm{clip}}.
\]

\subsection{clip fraction 的正确统计}

clip fraction 应只统计 response token：
\[
  f_{\mathrm{clip}}
  =
  \frac{\sum_{i,t}m_{i,t}\mathbf{1}[\text{token }(i,t)\text{ 被裁剪}]}
       {\sum_{i,t}m_{i,t}},
\]
其中 $m_{i,t}=0$ 的 prompt/padding token 不应进入分母。

\subsection{修复前后代码}

\begin{lstlisting}
# before
clipped_mask = clipped_pos | clipped_neg
clip_fraction = clipped_mask.float().mean()
\end{lstlisting}

\begin{lstlisting}
# after: cs336_alignment/grpo_helper.py
metadata = {
    "clip_fraction": clip_fraction,
    "clipped_mask": clipped_mask,
}
...
clipped_mask = metadata.pop("clipped_mask", None)
if clipped_mask is not None:
    metadata["clip_fraction"] = masked_mean(clipped_mask.float(), response_mask)
\end{lstlisting}

这样保留了 \texttt{compute\_grpo\_clip\_loss} 的局部元数据，又在 microbatch 聚合时使用 response mask 重新计算日志指标。

\section{Dr.GRPO：constant length normalization}

\subsection{问题}

GRPO 默认 token mean 损失为
\[
  \mathcal{L}_{\mathrm{token\_mean}}
  =
  \frac{\sum_{i,t}m_{i,t}\ell_{i,t}}
       {\sum_{i,t}m_{i,t}}.
\]
这会让长回答和短回答在 token 总量上被不同方式加权。Dr.GRPO 常用固定长度常数 $C$ 对每个样本归一化：
\[
  \mathcal{L}_{\mathrm{constant}}
  =
  \frac{1}{B}\sum_{i=1}^{B}
  \frac{\sum_{t=1}^{L}m_{i,t}\ell_{i,t}}{C}.
\]
若 $C$ 取最大生成长度，则每个样本的总贡献按同一尺度归一化。

\subsection{修复前代码}

\begin{lstlisting}
# before
policy_gradient_loss, metadata = compute_policy_gradient_loss(...)
loss = masked_mean(policy_gradient_loss, response_mask)
loss /= gradient_accumulation_steps
loss.backward()
\end{lstlisting}

\subsection{修复后代码}

\begin{lstlisting}
# after: cs336_alignment/grpo_helper.py
if loss_normalization == "token_mean":
    loss = masked_mean(policy_gradient_loss, response_mask)
elif loss_normalization == "constant":
    if normalize_constant is None:
        normalize_constant = policy_log_probs.shape[1]
    per_example_loss = masked_normalize(
        policy_gradient_loss,
        response_mask,
        normalize_constant=normalize_constant,
        dim=1,
    )
    loss = per_example_loss.mean()
else:
    raise ValueError(f"Unknown loss_normalization: {loss_normalization}")

loss /= gradient_accumulation_steps
loss.backward()
\end{lstlisting}

\subsection{adapter 层修复}

隐藏测试通常通过 \texttt{tests/adapters.py} 调用实现。修复前 adapter 未暴露 \texttt{loss\_normalization} 和 \texttt{normalize\_constant}，导致 Dr.GRPO 行为即使在 production helper 中实现，也无法被 adapter 路径测试。

\begin{lstlisting}
# after: tests/adapters.py
def run_grpo_microbatch_train_step(
    ...,
    loss_normalization: Literal["token_mean", "constant"] = "token_mean",
    normalize_constant: float | None = None,
):
    return grpo_microbatch_train_step(
        ...,
        loss_normalization,
        normalize_constant,
    )
\end{lstlisting}

\section{GRPO 训练循环：optimizer/scheduler 复用与真实 collate 路径}

\subsection{问题}

GRPO 每次 rollout 后都会对同一策略继续训练。如果 \texttt{train\_grpo\_model} 内部每次重新创建 optimizer/scheduler，则 AdamW 动量、学习率调度状态都会丢失：
\[
  m_t=\beta_1 m_{t-1}+(1-\beta_1)g_t,
  \quad
  v_t=\beta_2 v_{t-1}+(1-\beta_2)g_t.
\]
若每个 rollout 重置 optimizer，则 $m_{t-1},v_{t-1}$ 被清零，不再是连续优化过程。

\subsection{修复后代码}

\begin{lstlisting}
# after: grpo_experiment.py
def train_grpo_model(
    model: torch.nn.Module,
    tokenizer: "AutoTokenizer",
    grpo_dataset: Dataset,
    train_config: TrainConfig,
    global_step: int,
    optimizer: torch.optim.Optimizer,
    scheduler: torch.optim.lr_scheduler.LRScheduler,
):
    ...
\end{lstlisting}

optimizer/scheduler 在 \texttt{main()} 中创建一次，并传入每轮 rollout 的训练函数。新增 CPU smoke test 真实执行 \texttt{GRPODataset -> get\_grpo\_collate\_fn -> train\_grpo\_model}，同时捕获了缺失 \texttt{pad\_sequence} 的问题。

\subsection{pad\_sequence 修复}

\begin{lstlisting}
# before
old_log_probs = pad_sequence(old_log_probs_list, batch_first=True, padding_value=0.0)
# NameError: pad_sequence is not defined
\end{lstlisting}

\begin{lstlisting}
# after: cs336_alignment/grpo_experiment.py
from torch.nn.utils.rnn import pad_sequence
\end{lstlisting}

\section{无 GPU 环境与 vLLM 入口}

\subsection{问题}

原始脚本默认使用 CUDA：

\begin{lstlisting}
# before
train_device: str = "cuda"
eval_device: str = "cuda"
AutoModelForCausalLM.from_pretrained(
    train_config.model_path,
    torch_dtype=torch.bfloat16,
    attn_implementation="flash_attention_2",
)
llm = init__vllm(...)
\end{lstlisting}

在无 GPU 机器上，即使用户传入 \texttt{--train\_device cpu}，固定的 \texttt{flash\_attention\_2} 与 vLLM 初始化仍会失败。

\subsection{修复后加载参数}

\begin{lstlisting}
# after: experiment modules
def model_load_kwargs(train_device: str) -> dict:
    kwargs = {
        "use_cache": False,
        "trust_remote_code": True,
    }
    if str(train_device).startswith("cuda"):
        kwargs.update({
            "torch_dtype": torch.bfloat16,
            "attn_implementation": "flash_attention_2",
        })
    return kwargs
\end{lstlisting}

CPU 下不传 \texttt{flash\_attention\_2}，也不强制 \texttt{bfloat16}。

\subsection{SFT 的 CPU-only 行为}

SFT 训练本身不依赖 vLLM；vLLM 只用于评估。因此 SFT 在无 CUDA 时跳过 vLLM 初始化和评估：

\begin{lstlisting}
# after: sft_experiment.py
def can_use_vllm(device: str) -> bool:
    return str(device).startswith("cuda") and torch.cuda.is_available()

llm = None
if can_use_vllm(eval_config.eval_device):
    llm = init__vllm(...)
else:
    logger.warning("Skipping vLLM initialization/evaluation ...")
\end{lstlisting}

\subsection{EI / GRPO 的 CUDA 约束}

EI 和 GRPO 的核心 rollout 依赖 vLLM 生成，因此不能伪装成完整 CPU 可运行。修复后入口明确报错：

\begin{lstlisting}
# after: grpo_experiment.py / expert_iteration_experiment.py
def require_vllm_cuda(device: str):
    if not str(device).startswith("cuda") or not torch.cuda.is_available():
        raise RuntimeError("GRPO requires vLLM on CUDA for rollout generation/evaluation.")
\end{lstlisting}

这比在 vLLM/CUDA 深层栈中崩溃更可诊断，也符合当前机器无 GPU 的约束。

\section{数据路径硬编码}

\subsection{问题}

修复前数据加载逻辑依赖路径名：

\begin{lstlisting}
# before
if "gsm8k" in data_path:
    ...
if "MATH-500" in data_path:
    ...
\end{lstlisting}

如果相同 schema 的数据被放到其他目录，例如 \texttt{renamed.jsonl}，会加载出空的 prompts/cots/groundtruths。

\subsection{修复后代码}

\begin{lstlisting}
# after: cs336_alignment/data_utils.py
for entry in datasets:
    if {"problem", "solution", "answer"}.issubset(entry):
        groundtruths.append(extract_MATH500_groundtruth(entry["answer"]))
        cots.append(f"{entry['solution']} </think> <answer> {entry['answer']} </answer>")
        prompts.append(wrap_prompt(entry["problem"], prompt_template_path))
    elif {"question", "answer"}.issubset(entry):
        groundtruths.append(extract_gsm8k_groundtruth(entry["answer"]))
        cots.append(format_cot_to_think_answer(entry["answer"]))
        prompts.append(wrap_prompt(entry["question"], prompt_template_path))
    else:
        logger.error(f"Entry has unsupported schema keys: {sorted(entry.keys())}")
\end{lstlisting}

现在判断依据是 JSONL entry 的字段结构，而不是文件路径。

\section{测试策略与验证结果}

\subsection{新增回归测试覆盖}

新增 \texttt{tests/test\_regressions.py} 覆盖：

\begin{itemize}[leftmargin=2em]
  \item 大 logits 熵计算 finite；
  \item 零方差组内 reward 产生零 advantage；
  \item Dr.GRPO constant length normalization；
  \item adapter 层透传 Dr.GRPO constant normalization；
  \item GRPO clip fraction 只统计 response mask；
  \item SFT 训练循环真实迭代 dataloader；
  \item GRPO optimizer/scheduler 复用接口；
  \item GRPO collate + train loop CPU smoke test；
  \item CPU-safe model load kwargs；
  \item 数据加载按 schema 而非路径名识别。
\end{itemize}

\subsection{当前机器上的验证}

编译检查：

\begin{lstlisting}[language=bash]
python3 -m py_compile \
  cs336_alignment/sft_experiment.py \
  cs336_alignment/expert_iteration_experiment.py \
  cs336_alignment/grpo_experiment.py \
  cs336_alignment/sft_helper.py \
  cs336_alignment/grpo_helper.py \
  cs336_alignment/optional_helpers.py \
  cs336_alignment/data_utils.py \
  tests/adapters.py \
  tests/conftest.py \
  tests/test_regressions.py \
  tests/test_dpo.py \
  tests/test_data.py
\end{lstlisting}

CPU-only 测试：

\begin{lstlisting}[language=bash]
PYTHONDONTWRITEBYTECODE=1 python3 -m pytest -q \
  tests/test_regressions.py \
  tests/test_sft.py -k 'not tokenize and not get_response' \
  tests/test_grpo.py \
  tests/test_metrics.py \
  tests/test_dpo.py \
  tests/test_data.py
\end{lstlisting}

结果：

\[
  36\ \text{passed},\quad
  3\ \text{skipped},\quad
  2\ \text{deselected}.
\]

其中 3 个 skip 来自当前环境的 \texttt{transformers -> sklearn/scipy -> NumPy 2.2.6} ABI 不兼容；测试已经改为在该环境下跳过 transformers 相关路径，使纯 PyTorch 的 SFT/GRPO/metrics 回归测试仍可执行。

\section{剩余注意事项}

\begin{enumerate}[leftmargin=2em]
  \item GRPO 标准差口径当前保留 PyTorch 默认 sample std，以匹配课程 snapshot。若后续按论文公式严格采用 population std，应将 \texttt{std(..., correction=0)} 与 snapshot 一并更新。
  \item EI/GRPO 完整实验仍需要 CUDA + vLLM；当前无 GPU 机器只能验证 helper、collate、训练 microbatch 等 CPU 子路径。
  \item DPO 的数值测试依赖 transformers tiny model；当前机器因为 NumPy ABI 问题跳过，环境修复后应重新运行 \texttt{tests/test\_dpo.py}。
\end{enumerate}

\end{document}
